\section{The OntoDB Store}
\label{sec-store}

\sys is a single-process engine with the layers of a conventional disk-based
DBMS (Figure~\ref{fig-arch}).
A \sparql query is parsed, bound against the dictionary, planned as a
left-deep tree, reordered by the optimizer, and run by volcano-style
executors that pull rows of term ids.
The executors read triples through a store interface.
The \emph{indexed store}, which this paper is about, keeps each triple in
three B+ trees in one page file, behind a buffer pool with LRU-K
eviction~\cite{oneil1993}; terms are encoded by a dictionary built from the
same parts, a heap of spellings and two B+ trees.
Updates are logged page by page for ARIES-style recovery~\cite{mohan1992}.
The semantic extension of Sections~\ref{sec-labels} and~\ref{sec-vectors}
sits beside the store and reads it through the same interface.

\begin{figure}[t]
\centering
\begin{tikzpicture}[font=\scriptsize,
  box/.style={draw,rounded corners=1.5pt,minimum height=0.55cm,inner sep=3pt,align=center,fill=white,text width=3.1cm},
  side/.style={draw,rounded corners=1.5pt,minimum height=0.55cm,inner sep=3pt,align=center,fill=orange!12,text width=1.9cm},
  arr/.style={-{Stealth[length=1.6mm]},thick}]
\node[box] (q) {\sparql parser, binder, planner};
\node[box,below=0.22cm of q] (o) {optimizer (statistics), volcano executors};
\node[box,below=0.22cm of o,fill=blue!8] (s) {indexed store: SPO, POS, OSP};
\node[box,below=0.22cm of s,fill=blue!8] (t) {B+ tree: latch crabbing, sibling links};
\node[box,below=0.22cm of t,fill=blue!8] (b) {buffer pool (LRU-K), page guards};
\node[box,below=0.22cm of b] (d) {disk scheduler, disk manager};
\node[side,right=0.45cm of s] (dict) {dictionary\\heap + 2 B+ trees};
\node[side,right=0.45cm of b] (log) {log manager,\\recovery};
\node[side,left=0.45cm of o,text width=1.5cm] (sem) {semantic\\extension};
\draw[arr] (q) -- (o); \draw[arr] (o) -- (s); \draw[arr] (s) -- (t); \draw[arr] (t) -- (b); \draw[arr] (b) -- (d);
\draw[arr] (s) -- (dict); \draw[arr] (b) -- (log);
\draw[arr] (sem) |- (s);
\end{tikzpicture}
\caption{The layers of \sys. Blue boxes are the storage structures of Section~\ref{sec-store}.}
\label{fig-arch}
\end{figure}

\subsection{Keys, Pages, and Permutations}

A triple is a 24-byte key $x$ of three dictionary ids, stored big-endian so that
\texttt{memcmp} is the order.
A page is 4\,KiB; a leaf holds up to $B{=}169$ keys and an internal page up
to $F{=}144$ separators with $F{+}1$ child ids.
Leaves are linked to their \emph{right} sibling only.
A header page names the root, so a root split changes one word.
Every page but the root keeps at least half its maximum: a leaf at least
$\lfloor B/2\rfloor=84$ keys, an internal page at least
$\lfloor F/2\rfloor{+}1=73$ children.

Three permutations are enough.
A pattern binds some of $s$, $p$, $o$, and with the permutations SPO, POS,
and OSP every pattern has one in which its bound positions form a prefix
(Table~\ref{tab-perms}); it is then one range scan that returns only
matching triples.
Algorithm~\ref{alg-pattern} (Appendix~\ref{apx-alg}) picks the permutation with the longest bound
prefix.
The other three permutations would make no pattern cheaper
(Section~\ref{sec-exp}).

\begin{table}[t]
\caption{The permutation $\pi$ that answers each pattern, and the length $\ell(\pi)$ of the bound prefix.}
\label{tab-perms}
\small
\resizebox{\columnwidth}{!}{%
\begin{tabular}{@{}lccccccc@{}}
\toprule
Pattern & $(s,\cdot,\cdot)$ & $(\cdot,p,\cdot)$ & $(\cdot,\cdot,o)$ & $(s,p,\cdot)$ & $(\cdot,p,o)$ & $(s,\cdot,o)$ & $(s,p,o)$ \\
$\pi$, $\ell(\pi)$ & SPO, 1 & POS, 1 & OSP, 1 & SPO, 2 & POS, 2 & OSP, 2 & SPO, 3 \\
\bottomrule
\end{tabular}}
\end{table}


\subsection{Search}

A point lookup and the start of every scan use one descent.
\textsc{ChildFor} picks the child whose range holds $x$: the last separator
that is $\le x$, or child 0 if none is.
The descent \emph{crabs}: it latches the child before it releases the parent,
so a concurrent split cannot move $x$ away between two
steps~\cite{lehman1981,graefe2011}.

\begin{algorithm}[t]
\caption{\textsc{LeafFor}$(x)$ and \textsc{Contains}$(x)$}\label{alg-search}
\small
latch the header page (shared); $P \gets$ the root; \lIf{$P$ is nil}{\Return nil}
latch $P$ (shared); release the header\;
\While{$P$ is an internal page}{
  $P' \gets \textsc{ChildFor}(P,x)$ \tcp*{binary search over the separators}
  latch $P'$ (shared); release $P$; $P \gets P'$\;
}
\Return $P$\;
\medskip
\textsc{Contains}$(x)$: $P \gets \textsc{LeafFor}(x)$; \Return whether $x$ is in $P$ (binary search)\;
\end{algorithm}

A lookup reads $h$ pages.
Every non-root leaf holds at least 84 keys and every non-root internal page
has at least 73 children, so a tree of $n$ keys has at most $n/84$ leaves
and height $h \le 2 + \log_{73}(n/84)$: at most four levels for the 1.3
million triples of LUBM(10), and three when the tree is bulk-loaded (about
7,800 full leaves under 54 internal pages).

\subsection{Insert}

Insert is \emph{pessimistic} crabbing.
It latches the header and then each page on the way down in write mode, and
releases every latched ancestor, the header included, as soon as the page
just latched is \emph{safe}, meaning it cannot split.
What is still latched at the leaf is exactly what a split can reach, so a
split never needs a latch it does not hold, and latches are taken in one
order everywhere: the header, then top to bottom, siblings left to right.

\begin{algorithm}[t]
\caption{\textsc{Insert}$(x)$}\label{alg-insert}
\small
latch the header (exclusive); \lIf{the tree is empty}{make a leaf holding $x$ the root; \Return true}
$\mathit{path} \gets [\,]$; $P \gets$ the root\;
\Repeat{$P$ is a leaf}{
  latch $P$ (exclusive)\;
  \If{$P$ is not full}{release every latch in $\mathit{path}$ and the header \tcp*{$P$ is safe}}
  push $P$ on $\mathit{path}$; $P \gets \textsc{ChildFor}(P,x)$\;
}
$P \gets \mathit{path}.\text{pop}()$; \lIf{$x \in P$}{\Return false}
insert $x$ into $P$; \lIf{$|P| \le B$}{\Return true}
\tcp{split the leaf: the left half stays, the right half moves, its first key is \emph{copied} up}
$R \gets$ new leaf with the upper half of $P$; $R.\text{next} \gets P.\text{next}$; $P.\text{next} \gets R$; $\sigma \gets R.\text{keys}[0]$\;
\While{true}{
  \If{$\mathit{path}$ is empty}{make a new root over $(P, R)$ with separator $\sigma$; set it in the header; \Return true}
  $Q \gets \mathit{path}.\text{pop}()$; insert $(\sigma, R)$ into $Q$ after $P$\;
  \lIf{$Q$ has at most $F$ separators}{\Return true}
  \tcp{split the internal page: the middle separator \emph{moves} up and leaves both halves}
  $P \gets Q$; $R \gets$ new page with the upper half of $Q$; $\sigma \gets$ the middle separator of $Q$\;
}
\end{algorithm}

A leaf split leaves two leaves of $\lfloor (B{+}1)/2\rfloor$ and
$\lceil (B{+}1)/2\rceil$ keys, so the half-full invariant holds.
Splits propagate only while the parent is full, so with random keys a split
costs $O(1/B)$ page writes per insert on average.
A split or a merge is the only update that spans pages, and
Section~\ref{sec-store-log} makes it atomic for recovery.

\subsection{Delete}

Delete crabs the same way with a different notion of safe: a page above its
minimum (for the root, more than one key), because only then can it not
underflow.
After removing $x$, a leaf below $\lfloor B/2\rfloor$ keys repairs itself
against a sibling under the same parent, preferring the left one.

\begin{algorithm}[t]
\caption{\textsc{Remove}$(x)$}\label{alg-remove}
\small
descend as in \textsc{Insert}, releasing ancestors when the page just latched is safe (above its minimum)\;
$P \gets$ the leaf; \lIf{$x \notin P$}{\Return false}
remove $x$ from $P$; \lIf{$P$ is the root or $|P| \ge \lfloor B/2\rfloor$}{\Return true}
\While{the page $P$ on top of $\mathit{path}$ is under its minimum and has a parent $Q$}{
  $R \gets$ the left sibling of $P$ under $Q$ (the right one if $P$ is first); latch siblings left to right\;
  \eIf{the entries of $P$ and $R$ fit in one page}{
    \textbf{merge}: move the right page's entries into the left (a leaf keeps its sibling link; an internal page pulls the separator down); delete the right entry from $Q$\;
  }{
    \textbf{redistribute}: move one entry from $R$ to $P$ (a leaf moves a key; an internal page rotates through $Q$) and fix the separator in $Q$\;
  }
  $P \gets Q$\;
}
\lIf{the root has one child left}{make the child the root and free the old root}
\end{algorithm}

A leaf emptied by a merge is \emph{kept}, empty, with its sibling link still
pointing on, rather than freed.
An iterator may already hold its id (next subsection), and an empty page it
can step through is cheaper than a protocol that tells every iterator the
page is gone.
Space reclaimed this way is not reused until the tree is rebuilt; freed
internal pages are unreachable once their parent is updated and are
released.

\subsection{Range Scan}

The half-open range $[a,b)$ is the primitive the semantic extension needs
(Section~\ref{sec-labels}); a prefix scan is a range between the smallest and
the largest key with that prefix.
\textsc{BeginRange} descends to the leaf for $a$, finds the first key
$\ge a$ by binary search, and hands the leaf to an iterator.


The iterator \emph{copies} the keys of a leaf that lie below $b$, remembers
the right sibling's id, and releases the leaf.
It holds no latch and no pin between rows, so a consumer that stalls, such as
a join waiting on its other input or a client reading slowly, blocks no
writer and pins no frame.
Reading the next leaf is one page read and no descent.
The cost is the copy and a weaker guarantee: between two leaves a writer may
change the tree, so a scan sees each leaf as it was when it was read, and
the transaction layer's locks, not the latches, give a query its isolation.
The iterator checks that keys increase along the chain and that a run of
empty leaves does not cycle, and stops with an error otherwise.

\begin{lemma}[Cost of a range scan]\label{lem-scan}
A scan that returns $r$ keys from an insert-only or bulk-loaded tree of
height $h$ reads at most $h + \lceil r/\lfloor B/2\rfloor\rceil$ pages, and at
most $h + \lceil r/B\rceil$ in a bulk-loaded tree.
\end{lemma}
\begin{proof}
The descent ends at the first leaf, which is one of the $h$ pages.
Every leaf after it holds at least $\lfloor B/2\rfloor$ keys, all inside the
range except possibly the last, so at most $\lceil r/\lfloor
B/2\rfloor\rceil$ more leaves are read.
A bulk-loaded leaf is full, which replaces $\lfloor B/2\rfloor$ by $B$.
\end{proof}

\subsection{Bulk Load}

Loading a file inserts hundreds of thousands of triples into an empty store,
and inserting one at a time leaves leaves about half full.
\sys instead sorts the triples once for each permutation and builds the tree
bottom-up.


Each page is written once, in order, and no page is split.
A bulk-loaded tree is full, which Section~\ref{sec-exp} measures.

\subsection{Buffer Pool, Dictionary, and Query Processing}

\stitle{Buffer pool}.
Pages live in a fixed set of frames.
A \emph{page guard} pins a frame and takes its latch for the guard's
lifetime, and its destructor releases both, so no path leaves a page pinned.
A frame is evicted by LRU-K with $K{=}2$~\cite{oneil1993}, which ranks a
frame by the time of its second most recent use and so is not flooded by one
long scan.
Reads and writes go through a disk scheduler.

\stitle{Dictionary}.
Terms (IRIs, literals with a language or datatype, blank nodes) have one
canonical spelling and get a dense 64-bit id in the order they first appear.
The mapping is stored in pages: spellings in a heap, \emph{id $\to$ record}
in a B+ tree keyed (id, page, slot), and \emph{spelling $\to$ id} in a B+
tree keyed (hash, id) with a 64-bit FNV-1a hash; a lookup scans the hash
prefix and compares the candidates' spellings, so a collision costs a
comparison and not a wrong answer.
Decoding an id costs a tree descent and a heap read, about $0.7\,\mu$s, and
a query that returns thousands of rows decodes thousands of ids.
A direct-mapped cache of the 65,536 most recently decoded spellings (slot
$=$ id mod size) removes most of that on a warm store; ids are dense and a
spelling never changes, so an entry is never stale.
Section~\ref{sec-exp} says where the cache helps and where it does not.

\stitle{Planning and joins}.
The planner builds a left-deep tree of triple scans in textual order and the
optimizer reorders it greedily: it appends, at each step, a pattern that
shares a variable with what is already joined (any pattern, if none does),
the one whose estimated join output is smallest.
The output of a join of inputs of $a_1$ and $a_2$ rows is estimated as
$a_1a_2/\max(V_1,V_2)$, where $V_1$ and $V_2$ are the numbers of distinct values
of the join variable, from exact counts that the store
collects at load for every position, every pair of positions, and every
triple.
The join operator is a nested loop, an index nested loop (the inner scan is
reopened with the outer row's bindings as constants, so it is again a prefix
scan of one tree), or a hash join, chosen per node.
The \sparql subset is basic graph patterns with \texttt{FILTER},
\texttt{DISTINCT}, \texttt{ORDER BY}, \texttt{LIMIT}, \texttt{OFFSET}, and
\texttt{INSERT DATA} and \texttt{DELETE DATA}; \texttt{OPTIONAL},
\texttt{UNION}, property paths, and aggregates are not implemented.

\subsection{Transactions and Recovery}
\label{sec-store-log}

Updates run under strict two-phase locking with a deadlock detector.
Recovery is page-oriented: when a dirty page is released, the pool appends a
record holding the page's image and stamps the page with its log sequence
number, and before a page is written the pool forces the log to that stamp.
Redo therefore copies page images whose stamp is older than the record and
needs no knowledge of the B+ tree layout.
Undo of a triple is \emph{logical} and writes a compensation record, because
other transactions may have changed the page since.
A split or merge is the one update that spans pages: insert and remove open a
\emph{page group} that holds the log back from the group's first record
until it ends, and the pool does not choose a page of the group as a victim,
so a crash never leaves half a split on disk.
Commit forces the log, and checkpoints are sharp (flush every page, then
write the record).
A crash harness kills the process between the pages of every multi-page
operation and checks the tree afterwards: sorted keys, separators consistent
with leaves, a leaf chain that visits every key once, and the half-full
bound.
